\documentclass[11pt]{article}
\usepackage[margin=2.5cm]{geometry}
\usepackage{hyperref}

\begin{document}

\noindent\today

\vspace{1cm}
\noindent Dear Editor,

\vspace{0.5cm}
\noindent We submit the manuscript \textbf{``Remote SAMsing: From Segment Anything to Segment Everything''} for consideration as a research article in the \textit{International Journal of Applied Earth Observation and Geoinformation}.

SAM2, the most widely adopted foundation model for image segmentation, cannot segment a large remote sensing image on its own. A single pass with its default settings leaves 30--70\% of the scene without any mask, and tiling the image produces fragmented segments at every tile boundary with no mechanism to reconcile them. These two problems prevent SAM2 from being used as a practical tool for Object-Based Image Analysis. Remote SAMsing solves both: a multi-pass algorithm with black mask focusing and adaptive threshold decay raises coverage to 91--98\%, and a parameter-free boundary merge produces a single consistent label map for images of any size.

What sets this work apart from existing SAM-based remote sensing studies is the scope of the problem it addresses. Prior work focuses on improving per-object mask quality on individual patches. Remote SAMsing is, to our knowledge, the first pipeline that takes a large remote sensing image as input and produces a complete, spatially consistent segmentation as output, without requiring training data or architectural modifications to SAM2. The experimental evaluation spans seven scenes across three spatial resolutions (5~cm to 4.78~m GSD), two spectral compositions (natural RGB and MNF false-color), and two landscape types (urban and agricultural). A scalability test on a full Potsdam mosaic of $36{,}000 \times 54{,}000$ pixels (1.94 billion pixels) produced 124{,}180 segments at 97\% coverage with 81.8\% Achievable Segmentation Accuracy, confirming that the pipeline operates at production scale on a single consumer GPU.

The analysis also reveals a finding of practical value: tile size functions as an implicit scale parameter that controls a trade-off between coverage and per-object detection quality. Reducing tile size from 1{,}000 to 250 pixels raises detection from 56\% to 85\% on urban scenes with small objects, outperforming SAM2's own built-in multi-scale mechanism at comparable processing time. This gives practitioners a single, interpretable knob to adapt the pipeline to different scene types.

Remote SAMsing achieves Boundary IoU 3--8$\times$ higher than SLIC and Felzenszwalb baselines across all tested scenes, indicating that SAM2's learned boundaries align with objects far more precisely than spectral-based segmentation methods. The pipeline and all evaluation code are publicly available at \url{https://github.com/osmarluiz/sam-mosaic}.

This manuscript has not been published or submitted elsewhere. All authors have approved the manuscript and agree with its submission.

\vspace{1cm}
\noindent Sincerely,

\vspace{0.8cm}
\noindent Prof.\ Dr.\ Osmar Ab\'{i}lio de Carvalho J\'{u}nior\\
Department of Geography\\
University of Bras\'{i}lia, Brazil\\
\href{mailto:osmarjr@unb.br}{osmarjr@unb.br}

\vspace{0.3cm}
\noindent On behalf of all authors:\\
Osmar Luiz Ferreira de Carvalho, Anesmar Olino de Albuquerque, Daniel Guerreiro e Silva

\end{document}
